% 出席確認クイズ 第08回　データの収集と整備
%   attendance/attendance08.html から生成（解説は本ガイド用に加筆）。

\quizq{1}{「構造化データ」の例として適切なものはどれでしょうか?}%
  {音声ファイル}%
  {\correct{表形式の売上データ}}%
  {自由記述のアンケート回答}%
  {写真}%
  {正解は (b)。構造化データは行と列で整理された表形式のデータで、売上台帳や会員名簿が典型です。(a) 音声、(c) 自由記述、(d) 写真は非構造化データで、生成AIの普及によってこれらも分析しやすくなりました。}

\quizq{2}{「一次データ」と「二次データ」の説明として適切なものはどれでしょうか?}%
  {\correct{一次データは自ら収集したデータ、二次データは他者が収集・公開したデータ}}%
  {一次データは無料、二次データは有料}%
  {一次データは古く、二次データは新しい}%
  {一次データは数値、二次データは文字}%
  {正解は (a)。自らの目的のために収集したものが一次データ、政府や企業など他者が収集・公開したものを利用するのが二次データです。(b)(c)(d) の有料・新しさ・数値か文字かは、一次・二次の区別とは無関係です。二次データでは収集の目的と定義が自分の目的と合っているかの確認が必要です。}

\quizq{3}{データの前処理に含まれる作業として適切なものはどれでしょうか?}%
  {会議の設定}%
  {\correct{欠損値の処理や表記ゆれの統一}}%
  {商品の値付け}%
  {広告の作成}%
  {正解は (b)。前処理では、欠損値の処理、重複の削除、表記ゆれ（「株式会社」と「（株）」など）の統一、型の変換、外れ値の確認などを行い、分析できる状態に整えます。分析の作業時間の大半は前処理だと言われます。(a)(c)(d) はデータ整備の作業ではありません。}

\quizq{4}{日本の公的統計を横断的に入手できるポータルサイトはどれでしょうか?}%
  {\correct{e-Stat(政府統計の総合窓口)}}%
  {Googleマップ}%
  {e-Tax}%
  {マイナポータル}%
  {正解は (a)。e-Stat（政府統計の総合窓口）は、各府省が公表する政府統計を1つの窓口で検索・取得できるポータルサイトです。(c) e-Tax は国税の電子申告、(d) マイナポータルは行政手続きのサイトで、統計の入手先ではありません。}

\quizq{5}{生成AIツールに業務データを入力する際の注意として適切なものはどれでしょうか?}%
  {どんなデータでも自由に入力してよい}%
  {入力したデータは誰にも見られないので安全である}%
  {\correct{個人情報や機密情報を入力しないなど、組織の方針に従う}}%
  {個人情報を入力すると精度が上がるので推奨される}%
  {正解は (c)。生成AIツールに入力したデータは外部のサーバーに送られ、設定によっては学習に使われる可能性があるため、個人情報・機密情報の扱いは組織の方針に従います。(b)(d) は誤りで、入力したデータが「誰にも見られない」保証はなく、個人情報の入力で精度が上がるわけでもありません。}
